\documentclass[10pt]{article}
\usepackage[utf8]{inputenc}
\usepackage[T1]{fontenc}
\usepackage{lmodern}
\usepackage[margin=0.6in]{geometry}
\usepackage{enumitem}
\usepackage{titlesec}
\usepackage{xcolor}
\usepackage[hidelinks]{hyperref}
\definecolor{ink}{HTML}{1A2C54}
\definecolor{accent}{HTML}{7A1F2B}
\hypersetup{colorlinks=true, urlcolor=accent, linkcolor=accent}
\titleformat{\section}{\normalsize\bfseries\color{ink}}{}{0em}{}[{\color{accent}\titlerule[0.8pt]}]
\titlespacing{\section}{0pt}{8pt}{3pt}
\setlist[itemize]{leftmargin=1.1em, itemsep=0.5pt, topsep=1pt, parsep=0pt}
\setlength{\parindent}{0pt}
\pagestyle{empty}
\newcommand{\sep}{\;\textbar\;}

\begin{document}

\begin{center}
{\LARGE\bfseries\color{ink} Carlos César Chávez Padilla}\\[3pt]
\small \href{mailto:cesarchavezpadilla@gmail.com}{cesarchavezpadilla@gmail.com} \sep +1 (202) 696-5595 \sep
\href{https://cesarchavezp29.github.io}{cesarchavezp29.github.io} \sep
\href{https://scholar.google.com/citations?user=LsVe9yQAAAAJ}{Google Scholar}
\end{center}

\vspace{1pt}
{\small\itshape Economist working at the intersection of development, human capital, and public policy. Master in Public Policy (U.\ of Chicago) with hands-on experience at the IMF and the IDB, and builder of Qhawarina, a real-time open-data platform that nowcasts Peru's economy from web data.}

\section{Education}
\textbf{University of Chicago}, M.A.\ in Public Policy (Harris School), Merit Scholarship \hfill 2024--2025 \\
\textbf{Universidad Nacional Mayor de San Marcos}, B.A.\ Economics, honors, GPA 3.9/4.0 \hfill 2013--2017 \\
\textbf{Universidad Nacional de Ingeniería}, Diplomas in Applied Econometrics \& Finance \hfill 2018

\section{Policy \& Research Experience}
\textbf{International Monetary Fund} --- Research Analyst / Assistant, MCM Department \hfill 2020--2023
\begin{itemize}
\item Supported Article IV and country missions (Sweden, South Africa, Caribbean); produced macro-financial analysis used in surveillance work.
\item Co-authored IMF Working Paper 2022/180 on closing Peru's ethnic income gaps; quantified the policy levers behind inclusive growth.
\item Built a 180-country panel linking political instability to growth and productivity, informing risk assessment.
\end{itemize}
\textbf{Inter-American Development Bank} --- Research Consultant, Pension Fund Department \hfill 2019--2020
\begin{itemize}
\item Analyzed pension-system design and old-age coverage gaps across Latin America.
\end{itemize}
\textbf{University of Chicago} --- Research Professional, Center for the Economics of Human Development \hfill 2024--present
\begin{itemize}
\item Work with Prof.\ James J.\ Heckman on early-childhood policy and skill formation; translate structural results into policy-relevant terms.
\end{itemize}

\section{Flagship Project --- Qhawarina}
Founder and sole builder of \textbf{Qhawarina}, a real-time open-data platform for Peru: scrapes 42{,}000+ products daily across all 25 regions and nowcasts GDP and inflation from prices, filling the gap left by lagged official statistics. Open, transparent, and updated daily. \href{https://cesarchavezp29.github.io/projects/qhawarina/}{cesarchavezp29.github.io/projects/qhawarina}

\section{Policy-Relevant Research}
\begin{itemize}
\item \textbf{Minimum wage \& ethnic gaps in Peru} --- who gains and who is left out (\textit{J.\ Economics, Race \& Policy}; bunching estimator working paper).
\item \textbf{Mining and local welfare} --- effects of extractive activity on inequality, income, and poverty (\textit{Mineral Economics}).
\item \textbf{Political instability \& macro outcomes} --- coups, productivity, and sovereign risk across countries.
\item \textbf{Election integrity} --- a forensic test of fraud allegations in Peru's 2021 runoff (\textit{under review, Political Analysis}).
\end{itemize}

\section{Skills}
\textbf{Data \& programming:} Stata, R (advanced); Python, SQL, web scraping, machine learning (proficient). \quad
\textbf{Communication:} policy briefs, data visualization, public dashboards. \quad
\textbf{Languages:} Spanish (native), English (fluent).

\end{document}
